\subsection{伽罗瓦群}

定义 2.——设 N 是域 K 的一个伽罗瓦扩张。N 的全体 K-自同构组成的群称为 N 在 K 上的伽罗瓦群，记作 Gal(N/K)。

设 N 是 K 的一个有限伽罗瓦扩张。则 N 是 K 的有限可分拟伽罗瓦扩张。因此（V，第 32 页，命题 4 及 V，第 54 页，命题 3），Gal(N/K) 的阶等于 {[}N : K{]}。我们稍后将证明：若 N 是 K 的伽罗瓦扩张且 Gal(N/K) 是有限的，则 N 在 K 上次数有限（V，第 66 页，定理 3）。

设 \(\Omega\) 是 K 的一个代数闭扩张，设 A 是一个可分多项式在 \(\Omega\) 中的根集，该多项式记作 \(f \in \mathbf{K}[\mathbf{X}]\) 。则域 \(\mathbf{N} = \mathbf{K}(\mathbf{A})\) 是 K 的伽罗瓦扩张。显然，N 的每个 K-自同构都保持 A 稳定，且由于 A 在 K 上生成 N，映射 \(\sigma \mapsto \sigma | \mathbf{A}\) 是 \(\operatorname{Gal}(\mathbf{N} / \mathbf{K})\) 到一个子群 \(\Gamma\) （集 A 的对称群 \(\mathfrak{S}_{\mathbf{A}}\) 的）上的同构，该子群称为多项式 \(f\) 的伽罗瓦群。由注 3（V，第 55 页）可知，若 \(x\) 与 \(y\) 属于 A，则下列性质等价：

\begin{enumerate}
\def\labelenumi{\alph{enumi})}
\item
  \(x\) 与 \(y\) 在 \(\mathbf{K}\) 上共轭，
\item
  \(x\) 与 \(y\) 属于 \(\Gamma\) 作用下的同一轨道，
\item
  \(x\) 与 \(y\) 是 \(f\) 的同一个不可约因子的根。
\end{enumerate}

特别地，\(f\) 不可约当且仅当 \(\mathbf{A}\) 非空且 \(\Gamma\) 传递地作用于 \(\mathbf{A}\) 上。

例.——1) 设 K 的特征异于 2，N 是 K 的一个二次扩张。若 \(x \in \mathbf{N} - \mathbf{K}\) ，则我们有 \(\mathbf{N} = \mathbf{K}(x)\) ，且 \(x\) 在 K 上的极小多项式形如 \(f(\mathbf{X}) = \mathbf{X}^2 - a\mathbf{X} + b\) ，其中 \(a, b \in \mathbf{K}\) 。于是我们有 \(f(\mathbf{X}) = (\mathbf{X} - x)(\mathbf{X} - y)\) ，其中 \(y = a - x\) ，故 \(y\) 与 \(x\) 共轭；由于 \(f(\mathbf{X})\) 可分，扩张 N 是伽罗瓦的。群 \(\operatorname{Gal}(\mathbf{N}/\mathbf{K})\) 有两个元素，它们诱导集 \(\{x, y\}\) 的两个置换。

2) 设 \(f = \mathbf{X}^3 + \mathbf{X}^2 - 2\mathbf{X} - 1 \in \mathbf{Q}[\mathbf{X}]\) 。多项式 \(f\) 是不可约的；因为否则它将有一个根 \(x \in \mathbf{Q}\) ；把 \(x = a / b\) 写成 \(a, b \in \mathbf{Z}\) 且 \(a\) 与 \(b\) 互素的形式，我们将得到 \(a(a^2 + ab - 2b^2) = b^3\) 与 \(a^3 = b(b^2 + 2ab - a^2)\) ；但这蕴含 \(a\) 整除 \(b\) 且 \(b\) 整除 \(a\) ，从而 \(x = \pm 1\) ，这是不可能的。令 \(\xi = e^{2\pi i / 7} \in \mathbf{C}\) ；则多项式 \(f\) 以 \(\alpha = \xi + \xi^{-1}, \beta = \xi^2 + \xi^{-2}, \gamma = \xi^3 + \xi^{-3}\) 为根。我们有 \(\beta = \alpha^2 - 2\) 与 \(\gamma = \alpha^3 - 3\alpha\) ，故扩张 \(\mathbf{Q}(\alpha)\) 在 \(\mathbf{Q}\) 上是伽罗瓦的。\(\mathbf{Q}(\alpha)\) 在 \(\mathbf{Q}\) 上的伽罗瓦群是 3 阶循环群，由一个元素 \(\sigma\) 生成，使得 \(\sigma(\alpha) = \beta, \sigma(\beta) = \gamma, \sigma(\gamma) = \alpha.\)

\begin{enumerate}
\def\labelenumi{\arabic{enumi})}
\setcounter{enumi}{2}
\item
  设 K = Q 并取 \(f = X^{3} - 2\) 。利用整数分解为素因子之积（I，第 51 页），我们容易看出 2 不是 Q 中元素的立方。因此多项式 f 不可约，因为否则它在 Q 中将有一个根。设 \(A = \{x_{1}, x_{2}, x_{3}\}\) 为 f 在 \(\Omega\) 中的根集，\(\Gamma\) 为 f 的伽罗瓦群。它传递地作用于 A 上；故其阶可被 3 整除。另一方面，商 \(j = \frac{x_{2}}{x_{1}}\) 异于 1 且我们有 \(j^{3} = 1\) 。故 j 满足关系式 \(j^{2} + j + 1 = 0\) ；而多项式 \(T^{2} + T + 1 = \left(T + \frac{1}{2}\right)^{2} + \frac{3}{4}\) 在 Q 中没有根，这表明（例 1）\([Q(j): Q] = 2\) 。于是 \([N: Q]\) 可被 2 整除，从而 \(\Gamma\) 的阶可被 6 整除。由于 \(\Gamma\) 包含在 6 阶的群 \(S_{A}\) 之中，我们有 \(\Gamma = S_{A}\) .
\item
  设 K 的特征为 \(p \neq 0\) ，令 \(\mathrm{K}(\mathrm{T})\) 为有理分式域，\(U = T^{p} - T\) 。令 \(E = K[U]\) 且 \(F = K[T]\) ，则多项式 \(f(X) = X^{p} - X - U\) （\(E[X]\) 中的）在 F 中有根 \(T, T + 1, \ldots, T + p - 1\) 。设 \(\sigma\) 为 F 的使 \(\sigma(T) = T + 1\) 的 K-自同构。我们有 \(\sigma^{i}(T) = T + i\) 及 \(\sigma(U) = U\) 。群 \(G = \{1, \sigma, \ldots, \sigma^{p-1}\}\) 是 p 阶循环群，其不变量域包含 E；由于 \([F: E] \leqslant p\) ，戴德金定理（V，第 27 页，推论 2）蕴含 E 是 G 的不变量域且 \([F: E] = p\) 。于是多项式 f 在 \(E[X]\) 中不可约；E 的扩张 F 是伽罗瓦的，其伽罗瓦群 G 是 p 阶循环群，而群 \(\Gamma\) 是 T、\(T + 1, \ldots, T + p - 1\) 的循环置换所成的群。
\end{enumerate}

关于此例的推广，见 V，第 93 页，例 2。

\begin{enumerate}
\def\labelenumi{\arabic{enumi})}
\setcounter{enumi}{4}
\tightlist
\item
  设 \(\mathbf{F} = \mathbf{K}(\mathbf{X}_1, \ldots, \mathbf{X}_n)\) 是 \(n\) 个未定元 \(\mathbf{X}_1, \ldots, \mathbf{X}_n\) 在 \(\mathbf{K}\) 上的有理分式域。令
\end{enumerate}

\[
s _ {k} = \sum_ {1 \leqslant i _ {1} <   \dots <   i _ {k} \leqslant n} X _ {i _ {1}} \dots X _ {i _ {k}}
\]

其中 \(1 \leqslant k \leqslant n\) ，并令 \(\mathrm{E} = \mathrm{K}(s_1, \ldots, s_n)\) ；我们用 \(f(\mathrm{T})\) 表示多项式

\[
\mathrm{T} ^ {n} - s _ {1} \mathrm{T} ^ {n - 1} + \dots + (- 1) ^ {n} s _ {n}.
\]

我们有 \(f(\mathrm{T}) = \prod_{i=1}^{n} (\mathrm{T} - \mathrm{X}_i)\) ，故 \(\mathbf{F}\) 是可分多项式 \(f(\mathrm{T}) \in \mathrm{E}[\mathrm{T}]\) 的一个分裂域。进而，对每个置换 \(\sigma \in \mathfrak{S}_n\) ，存在唯一一个 \(\mathbf{K}\) -自同构 \(h_\sigma\) （\(\mathbf{F}\) 的），使得 \(h_\sigma(\mathrm{X}_i) = \mathrm{X}_{\sigma(i)}\) （对 \(1 \leqslant i \leqslant n\) ）；我们有 \(h_\sigma(s_k) = s_k\) （对 \(1 \leqslant k \leqslant n\) ），故 \(h_\sigma\) 是 \(\mathbf{F}\) 的一个 E-自同构。换言之，\(\mathbf{F}\) 是 \(\mathbf{E}\) 的伽罗瓦扩张，且在根集 \(\{\mathbf{X}_1, ..., \mathbf{X}_n\}\) （\(f(\mathrm{T})\) 的）上的限制定义了 \(\operatorname{Gal}(\mathrm{F}/\mathrm{E})\) 到群 \(\mathfrak{S}_n\) 上的一个同构。特别地，\(\mathbf{E}\) 由满足下式的有理分式 \(f\) 组成：

\[
f \left(\mathrm{X} _ {\sigma (1)}, \dots , \mathrm{X} _ {\sigma (n)}\right) = f \left(\mathrm{X} _ {1}, \dots , \mathrm{X} _ {n}\right)
\]

（对每个 \(\sigma \in \mathfrak{S}_n\) ；参看 IV，第 67 页，推论。）

\begin{enumerate}
\def\labelenumi{\arabic{enumi})}
\setcounter{enumi}{5}
\tightlist
\item
  设 \(f\) 是次数 \(> 0\) 的首一多项式，且 \(K\) 的特征 \(\neq 2\) 。在 \(A\) 上定义一个全序，记作 \(\leqslant\) ，并令 \(\delta(f) = \prod_{\alpha < \beta} (\beta - \alpha)\) ，
\end{enumerate}

\((\alpha, \beta) \in \mathbf{A} \times \mathbf{A}\) ，并对每个 \(\sigma \in \mathfrak{S}_{\mathbf{A}}\) 令 \(\delta_{\sigma}(f) = \prod_{\alpha < \beta} (\sigma(\beta) - \sigma(\alpha))\) 。我们有

\(\delta_{\sigma}(f) = \varepsilon(\sigma)\ \delta(f)\) ，其中 \(\varepsilon(\sigma)\) 是 \(\sigma\) 的符号（I，第 64 页），且 \(\delta(f) \neq 0\) 。对每个 \(\tau \in \operatorname{Gal}(\mathbf{N}/\mathbf{K})\) ，我们有 \(\tau(\delta(f)) = \delta_{\tau|_{\mathbf{A}}}(f)\) 。因此 \(\Gamma\) 包含在交错群 \(\mathfrak{A}_{\mathbf{A}}\) 中当且仅当 \(\delta(f) \in \mathbf{K}\) 。此外，\(\delta(f)^2 = \prod_{\alpha < \beta} (\beta - \alpha)^2 =\)

\(d(f)\) 是多项式 \(f\) 的判别式（IV，第 81 页）。因此 \(\Gamma \subset \mathfrak{A}_{\mathbf{A}}\) 当且仅当 \(d(f)\) 是 \(\mathbf{K}\) 中某个元素的平方。于是在例 2 中我们有 \(d(f) = 49 = 7^2\) ，而在例 3 中 \(d(f) = -108\) （IV，第 85 页）。

设 N 是 K 的一个伽罗瓦扩张，L 是 N 的一个在 K 上伽罗瓦的子扩张。N 的每个 K-自同构 \(\sigma\) 都诱导 L 的一个 K-自同构 \(\sigma_{L}\) （V，第 55 页，注 1）。因此映射 \(\sigma \mapsto \sigma_{L}\) 是 \(\operatorname{Gal}(N/K)\) 到 \(\operatorname{Gal}(L/K)\) 中的同态，称为限制同态。

命题 3.——Gal(N/K) 到 Gal(L/K) 的限制同态是满射。

更一般地，考虑 N 的两个子扩张 L 与 L'，以及 L 到 L' 上的一个 K-同构 u。选取 K 的一个把 N 作为子扩张包含在内的代数闭包 \(\Omega\) （V，第 23 页，定理 2）。存在 \(\Omega\) 的一个在 L 上与 u 一致的 K-自同构 v（V，第 52 页，命题 1），且由于 N 是 K 的拟伽罗瓦扩张，v 诱导 N 的一个 K-自同构 \(\sigma\) （V，第 55 页，注 1）。换言之，Gal(N/K) 中的元素 \(\sigma\) 在 L 上与 u 一致。

